% ===== roadmap of this week =====
\begin{frame}{이번 주차 개요}
  \begin{itemize}
    \item \kb{8.1 팀 프로젝트} --- 4주 파이프라인으로 정형 데이터에
      Week 2--7 모델을 적용. \kq{산출물은 높은 성능이 아니라
      ``성능 숫자를 정직하게 얻은 절차''.}
    \item \kb{8.2 Peer-Review} --- 이번엔 네가 심사위원. 체크리스트로
      두 팀을 검토하고, 장·절을 명시한 반박 가능한 질문을 쓴다.
    \item \kb{8.3 중간고사 대비 총정리} --- 유도 템플릿 3가지를 종이에
      재현하고, 장별 핵심 질문으로 자가진단.
  \end{itemize}
  \vskip0.6em
  \begin{block}{이번 장을 마치면}
    모델·지표·전처리·분할을 스스로 골라 근거지을 수 있고,
    3분할·베이스라인·``test 한 번'' 절차를 실제로 수행해
    보고서·발표로 제시할 수 있다.
  \end{block}
  \vskip0.4em
  \begin{center}
  \kq{같은 절차를 세 가지 언어(코드, 보고서, 리뷰)로 말할 수 있어야
       이 프로젝트가 끝난 것이다.}
  \end{center}
\end{frame}
